\label{sec:history}

\subsection{Pre-1941: The Hamilton Era and Its Problems}

From 1850 to 1900, Congress used the Hamilton method (largest remainders).
Hamilton's appeal was intuitive: give each state the integer part of its
proportional share, then distribute remaining seats to states with the
largest fractional parts.
However, the 1880 Census produced a devastating demonstration of Hamilton's
weakness: the \emph{Alabama Paradox}.

When the House was sized at 299 seats, Alabama received 8 seats.
When the House was increased to 300 seats, Alabama's seat count fell to 7
--- a counterintuitive result that no method based on monotone priority
could produce.
The paradox arose because Alabama's fractional remainder, which had ranked
it 8th among states competing for the last few seats at $H=299$,
fell below the threshold at $H=300$ when a new state displaced it in the
fractional ranking.

Congress responded by fixing the House size at each census to avoid
triggering the paradox.
The Apportionment Act of 1929 \citep{apportionment1929} temporarily froze
the House size at 435 and required automatic apportionment after each
census, but it did not specify which mathematical method to use.
This led to conflict: the 1930 Census was reapportioned by two different
methods (Hamilton and Webster) producing different results,
and Congress could not agree on which to certify.
The apportionment ultimately used for the 1930s was the same as
the 1910 apportionment --- no reapportionment occurred.

\subsection{The 1941 Act: Adoption of Huntington-Hill}

The crisis of the 1930s non-apportionment prompted Congress to resolve
the method question definitively.
The National Academy of Sciences convened a panel of mathematicians,
which in 1929 recommended Huntington-Hill (the method Edward V.
Huntington had been advocating since 1911 and which Hill~\citeyearpar{hill1911apportionment}
had independently proposed).

The Act of November~15, 1941 codified Huntington-Hill as the
statutory apportionment method (2~U.S.C.\ \S2a).
The Act's logic reflected three congressional concerns:
\begin{enumerate}
  \item \textbf{Paradox immunity}: Huntington-Hill, as a divisor method,
    cannot produce the Alabama Paradox.
    The House size can be adjusted without risking a state losing seats.
  \item \textbf{Mathematical neutrality}: The geometric mean rounding rule
    was characterised by Huntington as the most defensible on grounds of
    equal treatment of each seat within a state's allocation.
  \item \textbf{Near-proportionality}: The relative deviation criterion
    (Theorem~\ref{thm:hh-optimal}) provided a principled basis for
    Congress to defend the method in court.
\end{enumerate}
Webster was also seriously considered; the difference between the two
methods was primarily the treatment of small states near the rounding
boundary.
Congress's choice of HH over Webster has been described as reflecting
a subtle preference for protecting small-state representation
\citep{balinski1982fair,schmeckebier1941congressional}.

\subsection{Reapportionments Under 2 U.S.C.\ \S2a}

The 1941 Act has governed every apportionment since:
1950, 1960, 1970, 1980, 1990, 2000, 2010, and 2020.
In each case, the Census Bureau publishes its computation of the
priority sequence, identifies the last seat awarded and the state
that narrowly missed a seat, and certifies the final seat counts.
For the 2020 Census, the last seat (seat 435) went to New York;
Texas and Minnesota were within a small number of priority points
of gaining or losing a seat.
